% ==============================================================================
% APÊNDICE B
% ==============================================================================

\section{Reprodução do experimento}
\label{apd:apendice-b}

O código, a configuração e os resultados estão no repositório do projeto. Com o Docker em
execução, os comandos abaixo recriam o ambiente e reproduzem cada análise; todo o processamento
roda em contêineres com versões fixadas.

\begin{lstlisting}[numbers=none]
python tasks.py setup        # imagens, banco e Pagila do zero
python tasks.py smoke        # gera eventos conhecidos e confere o jsonlog
python tasks.py bench        # mede o custo do registro de log
python tasks.py simulate     # 28 dias simulados, semente 42
python tasks.py features     # janelas de 5 minutos, atributos e rótulos
python tasks.py evaluate     # ajuste na validação, teste, regras e sensibilidade
python tasks.py robustness --exp-id sim-s42-28d__w5m__ne200-ms1024-mf0.5-log1
python tasks.py report       # relatório de anomalias e figuras
python tasks.py confirm  --exp-id sim-s42-28d__w5m__ne200-ms1024-mf0.5-log1
python tasks.py ablation --exp-id sim-s42-28d__w5m__ne200-ms1024-mf0.5-log1
python tasks.py drift    --exp-id sim-s42-28d__w5m__ne200-ms1024-mf0.5-log1
python tasks.py test         # testes automatizados
\end{lstlisting}

Para usar o SACLA em um log real, sem rótulos, o comando \texttt{fit} aprende com um período
considerado normal, e o \texttt{score} analisa outro período e gera o relatório em HTML e CSV:

\begin{lstlisting}[numbers=none]
python tasks.py fit   --logs /logs --until 2026-03-16T03:00:00Z \
                      --model data/sacla/detector.joblib
python tasks.py score --logs /logs --since 2026-03-23T03:00:00Z \
                      --model data/sacla/detector.joblib --out data/sacla/relatorio
\end{lstlisting}
